
% PIGEAN VERSION 11182025
% Template Version: Version 3 October 2023
% See section 11 of the User Manual for version history
%
%%%%%%%%%%%%%%%%%%%%%%%%%%%%%%%%%%%%%%%%%%%%%%%%%%%%%%%%%%%%%%%%%%%%%%
%%                                                                 %%
%% Please do not use \input{...} to include other tex files.       %%
%% Submit your LaTeX manuscript as one .tex document.              %%
%%                                                                 %%
%% All additional figures and files should be attached             %%
%% separately and not embedded in the \TeX\ document itself.       %%
%%                                                                 %%
%%%%%%%%%%%%%%%%%%%%%%%%%%%%%%%%%%%%%%%%%%%%%%%%%%%%%%%%%%%%%%%%%%%%%

%%\documentclass[referee,sn-basic]{sn-jnl}% referee option is meant for double line spacing

%%=======================================================%%
%% to print line numbers in the margin use lineno option %%
%%=======================================================%%

%%\documentclass[lineno,sn-basic]{sn-jnl}% Basic Springer Nature Reference Style/Chemistry Reference Style

%%======================================================%%
%% to compile with pdflatex/xelatex use pdflatex option %%
%%======================================================%%

%%\documentclass[pdflatex,sn-basic]{sn-jnl}% Basic Springer Nature Reference Style/Chemistry Reference Style


%%Note: the following reference styles support Namedate and Numbered referencing. By default the style follows the most common style. To switch between the options you can add or remove “Numbered” in the optional parenthesis. 
%%The option is available for: sn-basic.bst, sn-vancouver.bst, sn-chicago.bst%  
 
\documentclass[sn-nature,lineno]{sn-jnl}% Style for submissions to Nature Portfolio journals
%%\documentclass[sn-basic]{sn-jnl}% Basic Springer Nature Reference Style/Chemistry Reference Style
%%\documentclass[sn-mathphys-num]{sn-jnl}% Math and Physical Sciences Numbered Reference Style 
%%\documentclass[sn-mathphys-ay]{sn-jnl}% Math and Physical Sciences Author Year Reference Style
%%\documentclass[sn-aps]{sn-jnl}% American Physical Society (APS) Reference Style
%%\documentclass[sn-vancouver,Numbered]{sn-jnl}% Vancouver Reference Style
%%\documentclass[sn-apa]{sn-jnl}% APA Reference Style 
%%\documentclass[sn-chicago]{sn-jnl}% Chicago-based Humanities Reference Style

%%%% Standard Packages
%%<additional latex packages if required can be included here>

% Changing geoemtry to look better (imo)
\geometry{a4paper, margin=1in}

% Change font to arial 
\usepackage{helvet}
\renewcommand{\familydefault}{\sfdefault}

\usepackage{graphicx}%
\usepackage{multirow}%
\usepackage{amsmath,amssymb,amsfonts}%
\usepackage{amsthm}%
\usepackage{mathrsfs}%
\usepackage[title]{appendix}%
\usepackage{xcolor}%
\usepackage{textcomp}%
\usepackage{manyfoot}%
\usepackage{booktabs}%
\usepackage{algorithm}%
\usepackage{algorithmicx}%
\usepackage{algpseudocode}%
\usepackage{listings}%
\usepackage{subcaption}
\usepackage{comment}
\usepackage{tikz}
\usepackage{luacode}
\usepackage{pdfpages}
\usepackage{hyperref}
%%%%

%%%%%=============================================================================%%%%
%%%%  Remarks: This template is provided to aid authors with the preparation
%%%%  of original research articles intended for submission to journals published 
%%%%  by Springer Nature. The guidance has been prepared in partnership with 
%%%%  production teams to conform to Springer Nature technical requirements. 
%%%%  Editorial and presentation requirements differ among journal portfolios and 
%%%%  research disciplines. You may find sections in this template are irrelevant 
%%%%  to your work and are empowered to omit any such section if allowed by the 
%%%%  journal you intend to submit to. The submission guidelines and policies 
%%%%  of the journal take precedence. A detailed User Manual is available in the 
%%%%  template package for technical guidance.
%%%%%=============================================================================%%%%

%% as per the requirement new theorem styles can be included as shown below
\theoremstyle{thmstyleone}%
\newtheorem{theorem}{Theorem}%  meant for continuous numbers
%%\newtheorem{theorem}{Theorem}[section]% meant for sectionwise numbers
%% optional argument [theorem] produces theorem numbering sequence instead of independent numbers for Proposition
\newtheorem{proposition}[theorem]{Proposition}% 
%%\newtheorem{proposition}{Proposition}% to get separate numbers for theorem and proposition etc.

\theoremstyle{thmstyletwo}%
\newtheorem{example}{Example}%
\newtheorem{remark}{Remark}%

\theoremstyle{thmstylethree}%
\newtheorem{definition}{Definition}%

%% New commands
% Probablity Dist function
\newcommand{\prob}{\text{Pr}}
% Variable lookup function
\newcommand{\getvar}[1]{\directlua{lookup("#1")}}

\raggedbottom
%%\unnumbered% uncomment this for unnumbered level heads

\begin{document}

\title[Article Title]{Beyond the Locus: Modeling Indirect Genetic Support to Expand the Reach of Human Genetics}

%Dual embeddings of 6,915 traits and 18,125 genes learned from direct and indirect human genetic support
%A dense statistical map of direct and indirect genetic support between 18,125 genes and 6,915 rare and common human traits


%%=============================================================%%
%% GivenName	-> \fnm{Joergen W.}
%% Particle	-> \spfx{van der} -> surname prefix
%% FamilyName	-> \sur{Ploeg}
%% Suffix	-> \sfx{IV}
%% \author*[1,2]{\fnm{Joergen W.} \spfx{van der} \sur{Ploeg} 
%%  \sfx{IV}}\email{iauthor@gmail.com}
%%=============================================================%%

\author*[1,2]{\fnm{First} \sur{Author}}\email{iauthor@gmail.com}

\author[2,3]{\fnm{Second} \sur{Author}}\email{iiauthor@gmail.com}
\equalcont{These authors contributed equally to this work.}

\author[1,2]{\fnm{Third} \sur{Author}}\email{iiiauthor@gmail.com}
\equalcont{These authors contributed equally to this work.}

\affil*[1]{\orgdiv{Department}, \orgname{Organization}, \orgaddress{\street{Street}, \city{City}, \postcode{100190}, \state{State}, \country{Country}}}

\affil[2]{\orgdiv{Department}, \orgname{Organization}, \orgaddress{\street{Street}, \city{City}, \postcode{10587}, \state{State}, \country{Country}}}

\affil[3]{\orgdiv{Department}, \orgname{Organization}, \orgaddress{\street{Street}, \city{City}, \postcode{610101}, \state{State}, \country{Country}}}

%%==================================%%
%% Sample for unstructured abstract %%
%%==================================%%

\abstract{
Human genetic evidence is known to increase the likelihood of clinical success in drug development \cite{ochoa_human_2022,rusina_genetic_2023}. Yet fewer than 5\% of active drug programs are supported by germline genetic data \cite{minikel_refining_2024}, revealing a critical gap in how genetic knowledge is incorporated into therapeutic pipelines. To help bridge this gap, we introduce the concept of indirect genetic support, which captures evidence derived from associations with biologically related genes linked through shared annotations. This expands on direct genetic support by probabilistically integrating both genetic and non-genetic information into a unified model aimed at improving mechanistic understanding and clinical outcomes. We present PIGEAN, a Bayesian framework that quantifies direct and indirect support using GWAS summary statistics, exome-wide gene-level associations, and rare disease gene sets. Applied across 6,000 complex traits, 2,000 rare diseases, 600,000 gene annotations, and 29,000 drug programs, PIGEAN constructs a weighted graph that serves as a genetically grounded backbone for biomedical discovery. Our analysis shows that drug targets prioritized via this graph achieve clinical success rates comparable to those found in curated drug target databases, while also revealing over 1,500 additional viable targets through indirect evidence alone. By organizing gene prioritization along the interpretable axes of direct and indirect support, PIGEAN offers a principled alternative to existing ad hoc methods. When paired with AI agents, this framework enables scalable, mechanistically informed exploration of the gene–phenotype landscape, probablistically connecting every gene to every trait using the structure of current biological knowledge.
}

%NOTES FOR ABSTRACT
%Novel inferences of course
%But beyond that, the presence of known and novel all in one place enables analyses (RS, gene function queries, gene lists for disease, gene lists for collection of phenotypes ) that could not be done previously
%In particular, shows shared basis of rare and common diseases, which increases further when considered at pathway level and not just genes


%%% Old abstract
% \abstract{Human genetic support increases the likelihood of therapeutic success and helps prioritize genes for further biological study. Here, we introduce the idea of indirect genetic support, which we distinguish from the traditional concept of direct support in that genetic associations for other genes support a target gene through shared gene annotations. We derive a mathemtical model, PIGEAN, that probabilistically estimates indirect genetic support from GWAS summary statisics, exome gene-level association statistics, and/or rare disease gene lists; we show that this model generalizes the idea of "bottom-up" or "convergence" based methods for gene prioritization and enables its combination with estimates of direct genetic support. We apply this model across \getvar{NUM_COMPLEX_TRAITS} complex traits, XXX rare diseases, and XXX gene annotations, assigning X\% of genes support for at least one trait. [Result from analysis of gene annotations]. [Result from analysis of how many GWAS loci have known mechanisms]. [Result from joint vs. marginal estimates and mechanisms]. [Result from drug target analysis]. These estimates of indirect support significantly expand the number of genes with links to human traits and enable probabilistic reasoning to replace ad hoc gene prioritization approaches that combine genetic associations with prior knowledge about gene function.
% }

%%================================%%
%% Sample for structured abstract %%
%%================================%%

% \abstract{\textbf{Purpose:} The abstract serves both as a general introduction to the topic and as a brief, non-technical summary of the main results and their implications. The abstract must not include subheadings (unless expressly permitted in the journal's Instructions to Authors), equations or citations. As a guide the abstract should not exceed 200 words. Most journals do not set a hard limit however authors are advised to check the author instructions for the journal they are submitting to.
% 
% \textbf{Methods:} The abstract serves both as a general introduction to the topic and as a brief, non-technical summary of the main results and their implications. The abstract must not include subheadings (unless expressly permitted in the journal's Instructions to Authors), equations or citations. As a guide the abstract should not exceed 200 words. Most journals do not set a hard limit however authors are advised to check the author instructions for the journal they are submitting to.
% 
% \textbf{Results:} The abstract serves both as a general introduction to the topic and as a brief, non-technical summary of the main results and their implications. The abstract must not include subheadings (unless expressly permitted in the journal's Instructions to Authors), equations or citations. As a guide the abstract should not exceed 200 words. Most journals do not set a hard limit however authors are advised to check the author instructions for the journal they are submitting to.
% 
% \textbf{Conclusion:} The abstract serves both as a general introduction to the topic and as a brief, non-technical summary of the main results and their implications. The abstract must not include subheadings (unless expressly permitted in the journal's Instructions to Authors), equations or citations. As a guide the abstract should not exceed 200 words. Most journals do not set a hard limit however authors are advised to check the author instructions for the journal they are submitting to.}

\keywords{keyword1, Keyword2, Keyword3, Keyword4}

%%\pacs[JEL Classification]{D8, H51}

%%\pacs[MSC Classification]{35A01, 65L10, 65L12, 65L20, 65L70}

\maketitle

\newpage
\tableofcontents
\newpage

\section{Main}\label{sec:introduction}

% \begin{figure}[h!]
%     \centering
%     {{canva:2}}
%     \caption{a. Not all genes that are trait relevant are associated with a genetic variant that rises to significance. b. Perhaps shared genetic annotations can allow us to infer trait relevance from genes that are missed by GWAS through a novel measure of indirect genetic support. c. If indirect genetic support can be put on the same scale as direct genetic support, we can recover relevant genes missed by GWAS.}
%     \label{fig:pigean-intuition}
% \end{figure}

% Broad introductory paragraph
Most human traits are influenced by many genes, and most human genes influence many traits. Mapping this complex web of gene–trait relationships is essential for uncovering disease biology \cite{cohen_sequence_2006}, understanding clinical heterogeneity \cite{udler_type_2018}, accurately diagnosing patients \cite{yang_clinical_2013}, and designing effective therapies \cite{plenge_validating_2013}. Among the most widely used and powerful approaches to to identify gene-trait relationships are genetic association studies. For common traits, genome-wide association studies (GWAS) have identified tens of thousands of trait-associated loci across thousands of phenotypes \cite{abdellaoui_15_2023}, while rare disease studies have linked thousands of Mendelian conditions to specific causal genes \cite{chung_meta-analysis_2023}. A genetic association provides a gene with ``genetic support" for a human trait -- establishing the human-relevance and causal effect of its perturbation.

The number of trait-associated genes identified by genetic association studies is, however, but a small fraction of the total number of trait-relevant genes -- that is, those that can be perturbed in humans to yield an effect on a trait of interest \cite{finan_druggable_2017}. Many trait-relevant genes escape detection (Figure \ref{fig:pigean-intuition-a}) due to limited statistical power \cite{wang_statistical_2019}, constraining evolutionary pressure on genetic variation \cite{gazal_linkage_2017}, or the difficulty of assigning noncoding GWAS signals to their effector genes  \cite{schipper_demystifying_2022,costanzo_realizing_2025}. Yet the full set of trait-relevant genes is of most interest for drug development and experimental study \cite{jhamb_pathway_2019,macnamara_network_2020}. In part due to this incompleteness, most translational research does not use human genetics to prioritize genes \cite{wehling_assessing_2009,dornbos_evaluating_2022}, even with the  increasing recognition that drug targets with genetic support are substantially more likely to succeed clinically \cite{plenge_validating_2013,minikel_refining_2024}. 

% Related Work and More specific stage setting to elicit the the contributions paragraph
By contrast, non-human-genetic approaches are more widely used to identify trait-relevant genes but come with critical tradeoffs. Experimental studies in animal \cite{seok_genomic_2013,widjaja_inhibiting_2019} and cell-based \cite{ben-david_genetic_2018,canon_clinical_2019} systems can establish causal links between genes and phenotypes, but their relevance to humans is often uncertain. Human observational studies, such as gene expression profiling \cite{wu_integrating_2022,rodriguez_machine_2021} or epidemiological associations \cite{barter_cholesteryl_2011,bhatt_cardiovascular_2019}, provide human-relevant context but lack causal directionality. The literature and biomedical databases are rife with gene-trait relationships derived from these or other non-genetic studies, and they are widely used as a starting point for drug discovery or experimental study \cite{swinney_how_2011,dornbos_evaluating_2022} (Figure \ref{fig:pigean-intuition-b}). Recent studies demonstrate how these sources can be rapidly and accurately queried with large language models (LLMs) \cite{toufiq_harnessing_2023,hu_evaluation_2025}. However, it is well documented that prioritizing genes using these evidence sources is prone to bias, incompleteness, or circularity \cite{brito_removing_2011,gillis_assessing_2013,gaudet_removing_2017,dunham_human_2018,haynes_gene_2018,stoeger_large_2018,richardson_meta_2024,tan_caution_2025}.

Ideally, a map of gene-trait relationships would combine the complementary benefits of human genetic and non-human-genetic evidence sources. It is commonplace to evaluate genes individually through manual expert synthesis of genetic data and the published literature \cite{mountjoy_open_2022,karjalainen_genome-wide_2024}, but a quantiative map across the genome and phenome requires scalabale statistical methods. Such methods have been developed both in the context of GWAS effector gene identification, where gene ``functional annotations'' enriched for genetic associations are used to prioritize genes nearby associated SNPs \cite{pers_biological_2015,weeks_leveraging_2023}, and in the context of rare disease diagnostics, where literature mining \cite{birgmeier_amelie_2020} and interaction networks \cite{smedley_next-generation_2015} are used to prioritize genes that carry variants observed in a patient. These approaches, however, rely on heuristics (such as harmonic averaging \cite{buniello_open_2025} or regressions using noisy or incomplete training data \cite{weeks_leveraging_2023,schipper_prioritizing_2025}) and often lack a biological model, which limits their portability across traits or application to settings other than that for which they were originally designed. As a result, the currently charted map of trait-gene relationships across rare and common diseases remains sparse, inaccurate, and inconsistent across traits.

% Contribution explanation
Here, we combine human genetic association data for {{num_common_traits}} common and {{num_rare_traits}} rare traits and non-human-genetic data encoded in {{num_available_gene_annotations}} gene annotations to systemaically identify trait-relevant genes (Figure 1A-C) to form a Genetic Effect Matrix over All Phenotypes (GEMAP). Each gene-trait relationship is a composite of two statistically calibrated values: a traditional measure of  ``direct genetic support'' (the strength of gene-trait association from human genetic data) and a new  measure of ``indirect genetic support'' (the extent to which the gene shares functional annotations with other trait-associated genes). Indirect genetic support is measured on the same scale as direct genetic support, estimates a measure of gene-trait relevance distinct from direct gene-trait association, and enables us to use unbiased human genetic associations to calibrate noisy and heterogeneous non-human-genetic data sources and learn which are trustworthy for each trait. We validated our method on three independent evidence axes—genetic, pharmacologic, and clinical. This map is orders of magnitude denser than manually curated knowledge graphs and can be used to calculate a novel resiliancy over every trait that provides novel insights into the mechansic heterogeneity of traits. 

%This map when thresholded is orders of magnitude denser than manually curated knowledge graphs and, when unthresholded, can be viewed as dual numerical embeddings of genes and traits learned from combined genetic and non-genetic data sources. We illustrate the utility of this resource to aid in the interpretation of gene function and trait mechanistic heterogeneity, predict drug target relative success more accurately and comprehensively than previous resources, and provide a resource for agentic AI not capturable by pretraining on text-only corpora.

%This decomposition formalizes and generalizes previous heuristic or machine learning-based integration strategies \cite{smedley_next-generation_2015,weeks_leveraging_2023,schipper_prioritizing_2025}, offering a statistically grounded alternative.

%Unlike conventional biomedical knowledge graphs \cite{li_graph_2022}, the PGEG avoids lossy edge labeling \cite{himmelstein_systematic_2017,morris_scalable_2023}, heuristic integration \cite{ochoa_next-generation_2023}, and manual curation \cite{putman_monarch_2024}, while remaining directly usable for downstream applications.


\section{Results}

% Results story
% 1. Show the embedding map of traits and how it overlays with ontological structures
% 2. Overview of the method
% 3. Show that indirect support is predictive of direct genetic support and controled for circularity
% 3. Show that indirect support is predictive of clinical success
% 4 (or 3.1). Show that both are really predictive of clinical success
% 5. Show which geneset libraries are good and bad for for which trait groups
% 6. Show how textmining model can be used to search literature in a genetics first way
% 7. Show how the embedding map can be used in agentic tooling for hypothesis generation and validation


\subsection{Calculating Indirect Genetic Support From Gene Annotations and Genetic Evidence}
\begin{figure}[h!]
    \centering
    {{canva:3}}
    \caption{Conceptual overview of the PIGEAN framework. (A) Demonstration of top-down prioritization methods from GWAS, WES and/or Mendelian disease databases. (B) Demonstration of bottom-up prioritization methods from gene annotations. (C) Graphical Model of PIGEAN with a depiction of fusing direct and indirect genetic evidence. (D) Positioning each gene across the continuum of direct and indirect support. (E) Leveraging this continuum to power downstream analysis pipelines which can be leveraged in agentic tooling (F) to yield AI-driven analysis at scale.}
    \label{fig:pigean-method}
\end{figure}

To infer the GEMAP across the entire phenome and genome, we sought to infer a probability of relevance for any gene to any trait. Doing so using available human genetic associations and non-human-genetic data posed several challenges. First, we needed to infer relationships across thousands of traits accurately, robustly, and efficiently. Second, the traits are heterogeneous in terms of genetic association data available (e.g. common variant associations for GWAS, Mendelian gene-disease validity assertions for rare diseases) and genetic architecture. Third, currently significant gene-trait associations cover only a small fraction of trait-relevant genes. Fourth, non-human-genetic datasets (i.e. curated pathway databases, results from animal model experiments, text-mined relationships, or gene expression enrichment analyses) are large in number and of varying quality, and it is unclear which datasets are useful for predicting trait-relevant genes. Finally, non-human-genetic datasets are biased toward well-studied genes and well-studied biological mechanisms.

To address these challenges, we developed a statistical model to quantify a gene's ``indirect genetic support'', and a method termed PIGEAN (Priors Inferred from GEne ANnotations) for homogeneously inferring this measure across common and rare traits. Indirect genetic support combines the unbiased rigor of genetic associations with the breadth and scope of gene annotations, which encode non-genetic-data either as gene sets (binary annotations) or gene quantitative vectors (continuous annotations) and can be constructed in a high-throughput manner from many biomedical data sources (\ref{sec:methods}). Indirect genetic support is an emergent property of a simple statistical model of observed genetic associations, in which gene annotations affect the probability that a gene is relevant to a trait, with effect sizes unknown and different for each annotation and each trait, and in which trait-relevant genes are more likely to produce genetic associations than are non-trait-relevant genes (\ref{sec:methods}).  This model statistically generalizes several widely used heuristic approaches for prioritizing genes from genetic association studies (\ref{sec:methods}), and it enables us to demonstrate mathematically that genetic associations not only directly implicate the genes for which they are observed but also indirectly implicate genes with shared annotations. PIGEAN applies Bayesian inference to this model to jointly infer gene annotation effect sizes and gene trait-relevance probabilities given a library of gene annotations and either GWAS summary statistics or rare disease gene lists, in so doing (a) using human genetic associations to select annotations with non-zero effects and calibrate their effect sizes and (b) using genetic annotations to identify additional trait-relevant genes beyond those with direct association evidence. Priors on the latent parameters enforce strict regularization that controls biases towards well-studied genes or profligate annotations, as well as enforcing model sparsity and interpretability (\ref{sec:methods}).

We studied PIGEAN across a universe of {{NUM_COMPLEX_TRAITS:,}} complex traits and {{NUM_RARE_DISEASE_GROUPS:,}} rare disease groups. For complex traits, we integrated full summary statistic files from {{NUM_GWAS_FULL_SUMSTATS}} GWAS covering {{NUM_TRAITS_FULL_SUMSTATS}} traits, as well as curated association lists from {{NUM_GWAS_CURATED}} GWAS across {{NUM_TRAITS_CURATED}} traits. We used an extension of a previously published approach for calculating a probablistic nearest gene method \cite{dornbos_evaluating_2022} to determine each gene's direct genetic support from variant-level associations (\ref{sec:methods}). For rare diseases, we mapped curated disease–gene confidence levels from Orphanet to probabilities of disease relevance (\ref{sec:methods}).  

To evaluate the contribution of non-genetic evidence, we constructed multiple PIGEAN models, each defined by a specific set of geneset libraries available during inference. These libraries collectively included {{NUM_TOTAL_GENE_ANNOTATIONS:,}} gene annotations spanning diverse sources (\ref{Supplementary Table 2}): knockout mouse phenotypes ({{NUM_MPO_TERMS:,}} Mammalian Phenotype Ontology terms), membership in pathways and curated gene sets ({{NUM_MSIGDB_GENESETS:,}} from the Molecular Signatures Database), co-occurrence in PubMedCentral manuscript figures ({{NUM_PFOCR_SETS:,}} sets from the Pathway Figure OCR project), co-occurrence with terms in PubMed abstracts ({{NUM_MESH_TERMS:,}} MeSH terms), network interaction partners ({{NUM_STRING_NEIGHBORS:,}} STRING neighbors), and single-cell gene expression patterns ({{NUM_SINGLE_CELL_CLUSTERS:,}} experimental clusters curated as part of the PoPS method).  

The resulting GEMAP (Figure \ref{fig:pigean-method}), generated from the top-performing model identified via model selection, produced a complete matrix of gene–trait relevance probabilities. This included {{NUM_GENETRAIT_REL_PP90:,}} gene-trait relationships ({{NUM_PER_TRAIT_PP90}} per trait and {{NUM_PER_GENE_PP90}} per gene) with a posterior probability (PP) above 90\% (i.e. false discovery rate < 0.1), {{NUM_GENETRAIT_REL_PP50:,}} gene–trait relationships ({{NUM_PER_TRAIT_PP50}} per trait and {{NUM_PER_GENE_PP50}} per gene) with PP>50\%, and {{NUM_GENESET_TRAIT_ASSOC_MILLIONS:.2f}}M significant (\ref{sec:methods}) gene set–trait associations. Each gene-trait relationship is supported by provenance that includes the specific SNP associations or curated gene records that produce direct support, and the specific gene set-trait relationships that produce indirect support. The full map is available for download, for interactive browsing via a web interface, or through AI-driven query to filter and explore results in response to high-level questions.  

\subsection{Indirect Support Interpretively Predicts Direct Genetic Support}

%1. Hold out cross validation to show indirect predicts direct
%2. Temporal to show combined predicts future
%3. Stratify across trait architectures and sample size, rare vs. common
%4. Bias towards well studied genes or traits with higher power
%5. Comparison to direct only baseline

% Figure 2 panels
%a. Show indirect predicts direct in hold one out chromosome (portal_model)
%b. Show mechanistic alignment in portal model results compared to null for combined
%c. Show relative success results compared to benchmarks for combined (portal model)
%d. Show it predicts future GWAS signal (portal) 

% Add how many more indirect genes vs only direct
% Plot - standard direct vs indirect plot and show orthongality.
% direct support is probablistic nearest gene
% Merge Cross chromosome lines for each trait
% Plot of sample size to demostrate calibration with direct support
% Another figure to show miami plot 

% What is a trait architecture?
% How do I show the bias towards well studied genes or traits with higher power?

% Think of threshold of cimbined as FDR (False Discovery Rate) based on Bayes Factors. TATTA paper has the FDR. 

% Relevant genes are if you perturb them they fuck up. However we can't perturn everything in humans to learn this because of ethics and cost.
% So how might relevant genes show up.
% - Successful drug targets would be a nice proxy, because they would target relevant genes no matter if they show up in GWAS.
% - Pertub seq data could also be leveraged maybe? That would support the causality of relevant gene definition. 


\begin{figure}[t!]
    \centering
    {{canva:4}}
    \caption{}
    \label{fig:validation-genetics}
\end{figure}

% Explanation of Panel A and B showing indirect predicts direct in rare and common disease.
Here, we evaluate indirect support along the genetic evidence axis spanning both common and rare traits. For common traits, we tested whether indirect support can predict direct genetic support using a leave-one-chromosome-out (LOCO) cross-validation across 12 validation traits (Supplementary Table \ref{tab:common-validation-traits}). In each fold, indirect support was learned from all chromosomes except one, and these scores were then used to predict direct genetic support on the held-out chromosome (Figure \ref{fig:validation-genetics}a). Indirect support significantly predicted direct genetic support for every held-out chromosome across all 12 traits (Supplementary Table \ref{tab:loco-validation}; trait-level $\beta$ range {{MIN_LOCO_BETA}} to {{MAX_LOCO_BETA}}, all $p$={{MAX_LOCO_BETA_P}}), and the fitted relationships consistently fell below the $y=x$ calibration line, indicating a systematic underestimation of direct support. This conservative bias is expected: NEED AN INTERPRETATION. Because LOCO validation is not feasible for rare diseases, we performed an analogous evaluation across 40 rare traits (Supplementary Table \ref{tab:rare-validation-traits}) using OMIM disease–gene annotations curated by Orphanet as a proxy for direct support. For each disease, we held out one known disease gene and asked how often it was recovered among the top-$K$ indirectly prioritized genes, summarized by a cumulative match characteristic (CMC) curve (Figure \ref{fig:validation-genetics}b). Indirect support recovered held-out rare disease genes with a CMC-AUC of {{RARE_CMC_AUC}} (95\% CI: {{RARE_CMC_AUC_CI_LOWER}} - {{RARE_CMC_AUC_CI_UPPER}}) and a Mean Reciprocal Rank (MRR) of {{RARE_MRR}} (95\% CI: {{RARE_MRR_CI_LOWER}} - {{RARE_MRR_CI_UPPER}}), where MRR is the average of $1/r$ across diseases (with $r$ being the rank of the held-out gene). An MRR of {{RARE_MRR}} implies the true gene is typically ranked within the top $~\sim$5-10 candidates (since $1/{{RARE_MRR}}\approx 8$) and far earlier than expected by chance in a genome-wide ranking. Consistent with this, {{RARE_RATIO_OF_PREDICTED_GENES_IN_TOP_1000}}\% of held-out genes appeared within the top 1000 predictions. In all analyses, we leveraged the most performant model chosen via model selection (Methods; Model Selection) that included two geneset libraries made up of mouse phenotypes and MSigDB genesets.

% Temporal validation of indirect support predicting future GWAS signal
Next, we sought to control for any potential circularity caused by genetic annotations that may have been informed by genetic data (GWAS or otherwise). We did this by curating a temporal validation set of 67 common traits that had a GWAS published before and after 2018 such that the GWAS published after 2018 had higher statistical power. We then leveraged the same Mouse + MSigDB model (described earlier) however using an older version of mouse phenptypes and MSigDB published before 2018. We then calculated gene-level MAGMA \cite{} Z-scores (as another independent metric of direct genetic support) from both the pre- and post-2018 GWAS for each trait. We then regressed the post-2018 MAGMA Z-scores on the pre-2018 MAGMA Z-scores for each trait while progressively restricting to genes meeting increasing thresholds of indirect genetic support (Figure \ref{fig:validation-genetics}c). We measured the improvement in model fit ($\Delta R^2$) of this regression as a function of the indirect genetic support threshold finding a statistically significant effect of {{TEMPORAL_DELTA_R2_BETA}} (p={{TEMPORAL_DELTA_R2_BETA_P}}) for indirect support to predict future GWAS signal. Together, these temporal validations support that PIGEAN is not simply recapitulating latent GWAS or drug target information embedded in curated libraries, but instead enriches for signals that are predictive of future discoveries. 

% conservative bias explanation, not all genes will have direct support on each axis but indirect support will still predict that they do even tho they won't rise to population significance level so we would expect the regression slope to be less than 1. Consistant with our expectation that not all disease relevant genes have variant association, we would expect this downward bias to be present. 
% Taking up what indirect support is, instead of validating it. 

% Show that the genesets used for prediction are interpretable and make sense and there is only a few that converge and that they also get deduplicated.
We then evaluated the interpretability of PIGEAN by examining the quantity and quality of significant genesets inferred by the same model over the entire set of traits (N={{NUM_TRAITS_MOUSE_MSIGDB_CONVERGED}}) as these effect sizes govern the calculation of indirect genetic support (Methods). We found that $\sim${{PERCENTAGE_OF_TRAITS_WITH_LESS_THAN_100_GENESETS}}\% ({{NUM_TRAITS_WITH_LESS_THAN_100_GENESETS}} of {{NUM_TRAITS_MOUSE_MSIGDB_CONVERGED}}) had fewer than 100 genesets predicted as significant (Figure~\ref{fig:validation-genetics}d), defined as joint effect size $\beta > 0.01$ (Methods). We then examined the quality of these significant genesets by calculating the semantic similarity between all significant genesets for a specific phenotype and the associated phenotype name and description (Figure~\ref{fig:validation-genetics}e). We found that the average phenotypic similarity (cosine similarity) was significantly greater than the random null ({{MEAN_PHENOTYPIC_SIMILARITY_OF_SIGNIFICANT_GENESETS}} compared to {{MEAN_PHENOTYPIC_SIMILARITY_OF_RANDOM_NULL}}; one-sided Welch's $t$-test with $t$={{T_STATISTIC_PHENOTYPIC_SIMILARITY}}, $p \ll 0.001$). Lastly, to validate how PIGEAN distributes signal across overlapping genesets, we investigated the relationship between marginal and joint geneset effects across traits. Under a simple approximation, marginal annotation effects can be viewed as joint effects filtered through the geneset overlap structure, so that overlapping genesets repeatedly capture the same latent biological signals, while joint effects deconfound this redundancy. Consistent with this redundancy-collapse interpretation, we observed that the total joint effect across genesets scales approximately with the square root of the total marginal effect (across $>${{NUM_TRAITS_MOUSE_MSIGDB_CONVERGED}} traits, $S_{\mathrm{joint}} \approx 0.8437\sqrt{S_{\mathrm{marg}}} - 1.02$; $R^2=0.769$; $p \ll 0.001$), suggesting that our geneset library redundantly encodes a smaller set of underlying biological processes that PIGEAN resolves into more interpretable joint effects.


\subsection{Indirect Support Prioritizes Drug Targets With Higher Relative Success}

\begin{figure}[t!]
    \centering
    {{canva:5}}
    \caption{Indirect and combined PIGEAN genetic support enrich for successful drug targets compared with Open Targets metrics.}
    \label{fig:validation-pharma}
\end{figure}

% Paragraph on continuous drug target relative success and comparing to Open Targets
We evaluated drug targets prioritized by PIGEAN's indirect genetic support, defining a target as a gene and semantically matched indication pair (Methods) with indirect support above a specified log-odds threshold (likelihood of relevance) across {{NUM_TOTAL_TRAITS_MOUSE_MSIGDB}} phenotypes, comprising {{NUM_COMMON_TRAITS_MOUSE_MSIGDB}} common traits and {{NUM_RARE_TRAITS_MOUSE_MSIGDB}} rare diseases. At the relative-success (RS) level achieved by Open Targets Genetics (OTG) with a minimum locus-to-gene (L2G) score $\geq 0.1$ as curated by Minikel et al.\ \cite{minikel_refining_2024} (RS = {{MINIKEL_OTG_RS:.2f}}; 95\% CI: {{MINIKEL_OTG_RS_CI_LOWER:.2f}}--{{MINIKEL_OTG_RS_CI_UPPER:.2f}}), PIGEAN's indirect support prioritized {{NUM_PIGEAN_PRIORITIZED_TARGETS_MOUSE_MSIGDB_AT_RS_eq_MINIKEL_OTG_INDIRECT}} targets, compared with only {{NUM_MINIKEL_OTG_PRIORITIZED_TARGETS}} targets prioritized by OTG. When direct and indirect genetic support were combined (combined genetic support) at the same RS level {{MINIKEL_OTG_RS}}, PIGEAN prioritized {{NUM_PIGEAN_PRIORITIZED_TARGETS_MOUSE_MSIGDB_AT_RS_eq_MINIKEL_OTG_COMBINED}} targets, but excluded {{NUM_UNIQUE_PIGEAN_INDIRECT_TARGETS_COMPARED_WITH_COMBINED}} targets uniquely identified by indirect support (with {{NUM_SHARED_PIGEAN_TARGETS_INDIRECT_AND_DIRECT}} targets shared between the indirect-only and combined sets), indicating that indirect and combined support surface partially distinct sets of drug targets.

Because OTG and its corresponding drug discovery resource, the Open Targets Platform (OTP), are the only settings in which a continuous RS metric is available, we benchmarked PIGEAN's indirect and combined genetic support against the Minikel et al.\ curation of OTG targets as a function of L2G score, as well as against OTP's harmonic mean score for genetic support and its literature- and model-based concept of indirect genetic support (Methods; Figure~\ref{fig:validation-pharma}a). RS was plotted as a function of the fraction of each resource prioritized at a given threshold of its associated score, parameterized by a threshold $\alpha \in [0,1]$. Across this range, PIGEAN's indirect and combined support consistently yielded higher RS than OTP's indirect-support and all-genetic-support metrics, while prioritizing a similar number of targets at $\alpha = 0.1$ (10\% of each resource). Although the Minikel et al.\ OTG curation performed slightly better than PIGEAN's indirect and combined support, this curated set was considerably smaller, containing only {{NUM_MINIKEL_OTG_TARGETS_AT_ALPHA_10}} targets at $\alpha = 0.1$.

We next examined OTG L2G score as a joint function of RS and the number of targets prioritized (Extended Data Figure~\ref{fig:extended-1}). L2G alone performed almost equivalently to PIGEAN's indirect genetic support, whereas combined PIGEAN support substantially outperformed all other approaches over the first 1,000 prioritized targets. Notably, RS for L2G remained relatively stable across the range of L2G scores, such that even scores close to 1 (nominally near-causal) did not yield markedly increased RS, consistent with the findings of Minikel et al. By contrast, both indirect and combined genetic support from PIGEAN achieved substantially higher RS at higher scores, indicating greater selectivity at more stringent thresholds.

% Paragraph on drug target success with temporal holdout of drug targets
Because the gene sets used by our PIGEAN models may contain latent information about drug targets, we further evaluated indirect and combined genetic support in a temporal holdout of drug targets over the previous 67 common traits with a GWAS published before 2018, using the 2018 versions of the MSigDB and mouse phenotype gene sets (same model as reported in Section 2.2). In this temporal analysis, we restricted evaluation to drug targets with launch dates after 2018. The Pharmaprojects dataset reports only the year of drug launch \cite{minikel_refining_2024} and does not capture the year a development program began, so we could exclude targets launched before 2018 but could not censor pre-2018 programs that did not launch; such programs were therefore retained in our RS calculations. Even under this temporal holdout, indirect and combined PIGEAN support remained strong predictors of RS for post-2018 drug targets (Figure~\ref{fig:validation-pharma}b), with indirect support achieving a maximum RS of {{MAX_RS_INDIRECT_TEMPORAL_HOLDOUT}} (95\% CI: {{MAX_RS_INDIRECT_TEMPORAL_HOLDOUT_CI_LOWER}}--{{MAX_RS_INDIRECT_TEMPORAL_HOLDOUT_CI_UPPER}}; $N$ = {{NUM_DRUG_TARGETS_SUPPORTED_INDIRECT_TEMPORAL_HOLDOUT}}) and combined support achieving a maximum RS of {{MAX_RS_COMBINED_TEMPORAL_HOLDOUT}} (95\% CI: {{MAX_RS_COMBINED_TEMPORAL_HOLDOUT_CI_LOWER}}--{{MAX_RS_COMBINED_TEMPORAL_HOLDOUT_CI_UPPER}}; $N$ = {{NUM_DRUG_TARGETS_SUPPORTED_COMBINED_TEMPORAL_HOLDOUT}}), with {{NUM_DRUG_TARGETS_SHARED_INDIRECT_AND_COMBINED_TEMPORAL_HOLDOUT}} targets shared between the indirect and combined sets. At $\alpha = 0.02$ (2\% of the resource), RS for indirect support stabilized at approximately {{RS_INDIRECT_TEMPORAL_HOLDOUT_STABLE_AT_ALPHA_002}} and RS for combined support at approximately {{RS_COMBINED_TEMPORAL_HOLDOUT_STABLE_AT_ALPHA_002}}, based on {{NUM_TARGETS_INDIRECT_TEMPORAL_HOLDOUT_AT_ALPHA_002}} and {{NUM_TARGETS_COMBINED_TEMPORAL_HOLDOUT_AT_ALPHA_002}} targets at this threshold, respectively.

% Paragraph on drug target success in comparision to minikel et al.

\subsection{Indirect Support Enrichs Genes for Rare Disease Patients}

Kyrung's Rare disease enrichments. 


\bibliography{references}% common bib file

% Commenting out for now because weird imports (need to fix)
% % Supplementary information
% \input{sections/indirect-support/01122026/supplementary_01122026}


\end{document}

% Some notes
% - pigean indirect + direct over common
% - pigean indirect over common (A2F and/or GWAS Cat)
% - combined
% - piccolo
% - OTG - All
% - magma 
% - bar or point
% - Temporal figure of indirect support regress magma2d
% - Take temporal direct support out 2c and cite in paper


% 3x2 for figure. Panel A is RS, Panel B is 1d

% We developed pigean and applied to XXXX common traits, predicted RS, and held when holding out future whatever. 
